\documentclass[12pt]{article}
%^ Start of document
\usepackage{times}  %Makes text TimesNewRoman
\usepackage{setspace} %Allows you to set single or double space 
\usepackage{biblatex} %Imports biblatex package
\usepackage{graphicx} % Required for inserting images
\usepackage{authblk}
\usepackage[colorlinks=true, urlcolor=blue, linkcolor=blue]{hyperref}
\usepackage{subfig}
\usepackage{float}
\usepackage{tabularx}
\usepackage{multirow}
\usepackage[paperheight=11in,
   paperwidth=8.5in,
   top=20mm,
   bottom=20mm,
   left=40mm,
   right=40mm]{geometry}
\usepackage{moresize}
\usepackage{tikz}
\usepackage{xcolor}
\usepackage{physics}
\usepackage{amssymb}
\usepackage{mathrsfs}
\usepackage{amsmath}
\usepackage{ragged2e}
\usepackage{comment}
\usepackage{xcolor}
\addbibresource{seminar.bib}
\begin{document}
\singlespacing  %Sets single spacing for whole document
\title{Neutron Stars - Alford}

\author[1]{Zachary Tullier}
\affil[1]{Student\\Physics Department\\ University of Houston - Clear Lake \\Houston, Texas\\U.S}
\date{}
\maketitle



\section{Introduction to the Speaker and Research}
\textbf{Speaker:} Mark Belford, a physicist based in St. Louis, specializes in studying super-dense matter within neutron stars. Originally from England, he has worked across various locations in the U.S. before settling in Missouri as a full-time physics professor.

\textbf{Research Focus:}
    Expertise in understanding quantum field theory applications related to the extreme conditions in neutron stars. With a focus in exploring matter states like superfluid quarks, dense nuclear matter, and the interplay of quantum mechanics with astrophysical phenomena.

\section{Overview of Neutron Stars}
    Neutron stars are remnants of massive stars that collapse after exhausting their nuclear fuel. They represent the densest form of observable matter in the universe, second only to black holes. They are typically 15-20 kilometers in diameter, with a mass comparable to or greater than the Sun, and densities billions of times higher than Earth materials.
    Their outer layer is composed of extremely dense normal matter. The inner layers gradually transition into a "neutron-rich soup," where protons and electrons combine to form neutrons.
    The core's composition is speculative and may involve exotic forms of matter like quarks or superfluid states.


\section{Formation of Neutron Stars}
    Stars begin as clouds of dust and gas that coalesce under gravity, forming a central core where nuclear fusion occurs. Depending on their mass, stars follow one of two life paths. One is small stars, similar to the Sun. They end as white dwarfs after passing through red giant and nebula phases. The other path is massive stars which are 20-30 times the Sun’s mass. They explode/collapse as supernovae, leaving behind either a neutron star or black hole.

\textbf{Nuclear Fusion and Iron Formation:}
    Stars generate energy by fusing lighter elements (hydrogen to helium, helium to carbon, etc.). The fusion process halts at iron, as iron cannot undergo further fusion to produce energy. This leads to core collapse and the creation of a neutron star.

\section{Detailed Structure and Composition}
    The outer crust is initially composed of compressed atomic matter, far denser than materials found on Earth. Further in is the transition zone. As pressure increases, protons in atomic nuclei absorb electrons, forming neutrons and releasing neutrinos. This process causes atoms to collapse into a denser configuration dominated by neutrons. Further in is the inner core. At the highest densities, matter exists as a uniform liquid of neutrons, protons, and electrons. The extreme environment may also lead to exotic phases, such as ``nuclear pasta'' (structured forms like rods or sheets of nucleons) or even a pure quark liquid.


\section{Exotic Matter and Theoretical Limits}
    Superfluid Quarks may exist in the interior. At extreme densities, neutrons may overlap to form a quark liquid, where individual quarks are no longer confined within protons and neutrons. This remains speculative but is a significant focus of current research. The degeneracy pressure from neutrons (a type of Fermi pressure) opposes gravitational collapse. However, if gravity overwhelms this pressure, the star collapses further into a black hole. Observations suggest neutron stars cannot exceed $\sim$2.1 solar masses. Beyond this, even degeneracy pressure cannot counteract gravity, leading to total collapse. Adding mass to a neutron star reduces its size due to the compressibility of nuclear matter under intense gravity. A typical neutron star ($\sim$1.4 solar masses) has a radius of about 10-14 kilometers. The mass-radius relation provides insights into nuclear forces at extremely high densities. Observations are compared with theoretical models to refine our understanding of these forces. Determining the radius of distant neutron stars is complex due to their small size and vast distances, leading to significant observational uncertainties.


\section{Neutron Stars and the Origin of Heavy Elements}
    The universe began predominantly with hydrogen and helium, with heavier elements forming later through stellar processes. Supernovae contribute to creating elements up to iron and slightly beyond. Heavier elements like gold, silver, and uranium are primarily produced during neutron star collisions. These events eject material into the surrounding space, enriching subsequent generations of stars and planets with heavy elements.


\section{Observational and Theoretical Tools}
    Measurements of neutron star masses and radii are critical for understanding their equation of state (pressure-density relation). Current data is imprecise, with significant error margins that limit our ability to form definitive conclusions. Researchers use quantum mechanical calculations to predict the properties of dense matter, such as viscosity, pressure, and phase transitions. These predictions feed into simulations of neutron star mergers, which are then compared to astronomical observations. The Tolman-Oppenheimer-Volkoff (TOV) equation describes how mass and radius relate for spherical neutron stars. Adjustments are required for spinning stars or systems under tidal forces.

\section{Challenges and Unresolved Questions}
    The behavior of nuclear forces at densities beyond twice that of atomic nuclei remains poorly understood. The transition from nucleons (protons and neutrons) to quarks is a key area of study, as this could dramatically alter the equation of state. Observational techniques and theoretical models are still evolving. Improved telescopes and computational methods are expected to narrow uncertainties in the future. Understanding the precise mass at which a neutron star collapses into a black hole could shed light on the ultimate limits of matter under gravity.


\section{Future Prospects}
    Advances in X-ray and gravitational wave astronomy hold the potential to provide more accurate measurements of neutron star properties. Better understanding of quantum chromodynamics (QCD), the theory governing quark interactions, could refine predictions for dense matter behavior. Insights from neutron star studies have implications for fundamental physics, astrophysics, and even nuclear energy technologies.

\section{Conclusion}
Neutron stars are natural laboratories for studying extreme physics, bridging astrophysics, nuclear physics, and quantum mechanics. Through a combination of observation, theory, and simulation, scientists aim to unravel the mysteries of dense matter, element formation, and the ultimate fate of collapsing stars. This research not only deepens our understanding of the cosmos but also challenges the limits of our knowledge about the fundamental nature of matter.

\end{document}



    
